% Part: computability
% Chapter: computability-theory
% Section: no-universal-function

\documentclass[../../../include/open-logic-section]{subfiles}

\begin{document}

\olfileid[tr]{cmp}{thy}{nou}
\olsection{Evrensel Hesaplanabilir Fonksiyon Yoktur}

Kısmi hesaplanabilir fonksiyonlar için evrensel olan kısmi hesaplanabilir bir
fonksiyon bulunsa da tümel hesaplanabilir fonksiyonlar için evrensel olan
tümel hesaplanabilir bir fonksiyon yoktur.

\begin{thm}
\ollabel{thm:no-univ}
Evrensel bir hesaplanabilir fonksiyon yoktur. Başka bir deyişle, $f(x)$
herhangi bir tümel hesaplanabilir fonksiyon olduğunda her $x$ için
$f(x)=\fn{Un}'(k,x)$ olacak bir $k$ doğal sayısının bulunmasını sağlayan
her $\fn{Un}'(k,x)$ fonksiyonu hesaplanabilir değildir.
\end{thm}

\begin{proof}
İspat basit bir köşegenleştirmedir. $\fn{Un}'(k,x)$ tümel ve hesaplanabilir
olsaydı
\[
  d(x)=\fn{Un}'(x,x)+1
\]
de tümel ve hesaplanabilir olurdu. Fakat tanım gereği
$d(k)\neq\fn{Un}'(k,k)$ olur. Dolayısıyla her $k$ için $d(x)$ ile
$\fn{Un}'(k,x)$ fonksiyonlarının değerleri en az bir $x$'te, özellikle
$x=k$ iken farklıdır.
\end{proof}

\begin{explain}
Yukarıdaki \olref[uni]{thm:univ-comp}, bu köşegenleştirme savının ancak
evrensel fonksiyonun kısmi olmasına izin verme pahasına aşılabildiğini
gösterir. Yani $\fn{Un}$ tümel hesaplanabilir fonksiyonlar için evrenseldir,
fakat kendisi tümel değildir. Köşegenleştirme savı kısmi durumda işlemez.
\end{explain}

\begin{prob}
  \olref{thm:no-univ} ispatındaki köşegenleştirme savının kısmi durumda neden
  işlemediğini anlamak için $f(x)\simeq\fn{Un}(x,x)+1$ fonksiyonunu ele alın.
  Bu fonksiyon kısmi hesaplanabilir midir? Öyleyse bir $e$ indisi vardır,
  yani $f(x)\simeq\fn{Un}(e,x)$ olur. $f(e)$ hakkında ne söyleyebilirsiniz?
\end{prob}

\end{document}

